\chapter{Vector Spaces}\label{ch:b1:vspaces}

Linear algebra begins here: the axioms of vector spaces isolate what
$\R^2$, $\R^3$, polynomial spaces and function spaces have in common
--- one can add, and scale. Two chapters build the theory
(\cref{ch:b1:findim} adds dimension); the language they set up ---
span, free family, basis, direct sum --- is the daily bread of every
chapter after them. Throughout, $K$ denotes $\R$ or $\C$
(\emph{scalars}).

\section{Definition and examples}

\begin{definition}[Vector space]\label{def:b1:vspaces:def}
A \emph{$K$-vector space}\index{vector space} is a set $E$ with an
addition making $(E, +)$ an abelian group (zero written $0_E$ or
$0$), and a scalar multiplication $K \times E \to E$ such that, for
all $\lambda, \mu \in K$ and $x, y \in E$:
\[
\lambda(x + y) = \lambda x + \lambda y,\quad
(\lambda + \mu) x = \lambda x + \mu x,\quad
\lambda(\mu x) = (\lambda\mu) x,\quad
1\,x = x .
\]
Consequences: $0\,x = 0_E$, $\lambda\,0_E = 0_E$, $(-1)x = -x$, and
$\lambda x = 0_E \implies \lambda = 0$ or $x = 0_E$ (multiply by
$\lambda^{-1}$).
\end{definition}

\begin{proof}[Proof of the consequences]
For $0\,x = 0_E$: from $(0 + 0)x = 0x + 0x$ and $(0+0)x = 0x$,
cancel $0x$ in the group $(E, +)$. For $\lambda\,0_E$: same
trick on $\lambda(0_E + 0_E)$. For $(-1)x$: add $x$,
\[
x + (-1)x = 1\,x + (-1)x = \bigl(1 + (-1)\bigr)x = 0\,x = 0_E ,
\]
so $(-1)x$ is the additive inverse of $x$. Finally if $\lambda x
= 0_E$ with $\lambda \neq 0$: multiply by $\lambda^{-1}$ (the
scalars form a field) and use the two axioms
$\lambda^{-1}(\lambda x) = (\lambda^{-1}\lambda)x = 1x = x$
together with $\lambda^{-1}0_E = 0_E$: $x = 0_E$. Small as they
are, these four rules are used silently on every page that
follows --- and the last one is exactly where fields are needed:
over the scalars $\Z$, the ``space'' $\Z/2\Z$ would violate it
with $2\,x = 0$.
\end{proof}

\begin{example}\label{ex:b1:vspaces:examples}
$K^n$ (coordinatewise operations); the polynomials $K[X]$; the
functions $\mathcal{F}(A, K)$ from any set $A$ to $K$ (pointwise
operations) --- containing continuous functions, sequences
$\mathcal{F}(\N, \R)$, etc.; $\C$ as an $\R$-vector space. In each
case the axioms are inherited from those of $K$.
\end{example}

\begin{definition}[Subspace]\label{def:b1:vspaces:subspace}
$F \subseteq E$ is a \emph{subspace}\index{subspace} when $0_E \in
F$ and $F$ is stable under addition and scalar multiplication ---
equivalently:
\[
F \neq \emptyset
\qquad\text{and}\qquad
\forall x, y \in F,\ \forall \lambda \in K,\quad
x + \lambda y \in F .
\]
A subspace is itself a vector space. Any intersection of subspaces
is a subspace; a union almost never is (same proof as
\cref{exo:b1:structures:6}).
\end{definition}

\begin{example}\label{ex:b1:vspaces:subspaces}
In $\mathcal{F}(\R, \R)$: the continuous functions, the
differentiable ones, the polynomials of degree $\leq n$ (written
$K_n[X]$ inside $K[X]$), the solutions of a homogeneous linear
differential equation (\cref{thm:b1:diffeq:homogeneous2} said just
that). Non-examples: $\{f : f(0) = 1\}$ (no zero), degree exactly
$n$ (not stable under addition).
\end{example}

\begin{example}[Subspace or not: four verdicts, argued]
\label{ex:b1:vspaces:subspaceornot}
In the space of real sequences:
\begin{itemize}
  \item $\{u : u \text{ bounded}\}$ \emph{is} a subspace: $0$ is
        bounded, and if $\abs{u_n} \leq M$, $\abs{v_n} \leq M'$,
        then $\abs{u_n + \lambda v_n} \leq M + \abs\lambda M'$.
  \item $\{u : u_n \to 1\}$ is \emph{not}: the zero sequence is
        missing (and the sum of two members tends to $2$).
  \item $\{u : u \text{ monotone}\}$ is \emph{not}: $u_n = n$
        and $v_n = -n + (-1)^n$ are monotone, their sum
        $(-1)^n$ is not; stability under addition is the axiom
        that fails, even though the set contains $0$ and all
        scalar multiples of its members.
  \item $\{u : u_{n+1} = u_n^2\}$ is \emph{not}: it contains $0$
        but $2u$ escapes as soon as $u$ is a nonzero member
        ($2u_{n+1} \neq (2u_n)^2$ in general) --- squaring is
        the nonlinearity.
\end{itemize}
The working order is always the same: test $0$ first (cheapest),
then stability --- and to refute, one explicit counterexample
pair beats any amount of doubt.
\end{example}

\section{Span, sums, direct sums}

\begin{definition}[Linear combinations, span]\label{def:b1:vspaces:span}
A \emph{linear combination} of the family $(x_1, \dots, x_p)$ of
vectors of $E$ is any $\lambda_1 x_1 + \dots + \lambda_p x_p$
($\lambda_i \in K$). The set of all of them is the
\emph{span}\index{span} $\operatorname{Vect}(x_1, \dots, x_p)$: it
is a subspace, the smallest one containing the family.
\end{definition}

\begin{proof}[Proof of the two assertions]
Stability: a sum of two linear combinations $\sum\lambda_i x_i +
\sum\mu_i x_i = \sum(\lambda_i + \mu_i)x_i$ is again one, and so
is a scalar multiple $\mu\sum\lambda_i x_i = \sum(\mu\lambda_i)
x_i$; the zero combination shows $0$ belongs: the span is a
subspace. Minimality: let $H$ be any subspace containing $x_1,
\dots, x_p$. By stability under scalar multiplication each
$\lambda_i x_i \in H$, and by stability under addition their sum
lies in $H$: every linear combination belongs to $H$, i.e.\
$\operatorname{Vect}(x_1, \dots, x_p) \subseteq H$. So the span
is contained in every subspace containing the family: it is the
smallest one.
\end{proof}

\begin{definition}[Sum, direct sum]\label{def:b1:vspaces:sum}
For subspaces $F, G$ of $E$:
\[
F + G = \{\,u + v : u \in F,\ v \in G\,\}
\]
is a subspace (the smallest containing $F \cup G$). The sum is
\emph{direct}\index{direct sum}, written $F \oplus G$, when every
element of $F + G$ decomposes \emph{uniquely} as $u + v$;
equivalently (see below) when $F \cap G = \{0\}$. When $E = F \oplus
G$, the subspaces are \emph{supplementary}\index{supplementary
subspaces} in $E$.
\end{definition}

\begin{example}[A sum of two lines]
\label{ex:b1:vspaces:sumoflines}
In $\R^3$, let $F = \operatorname{Vect}\bigl((1,0,1)\bigr)$ and
$G = \operatorname{Vect}\bigl((0,1,1)\bigr)$. Their sum is
\[
F + G = \{\,a(1,0,1) + b(0,1,1)\,\}
= \{(a,\ b,\ a + b)\} = \{(x, y, z) : z = x + y\},
\]
the plane through the origin containing both lines. It is
strictly bigger than the union $F \cup G$ (the mere cross of the
two lines): the vector $(1, 1, 2) = (1,0,1) + (0,1,1)$ lies in
the sum but on neither line. And $F \cap G = \{0\}$ (a common
vector requires $a(1,0,1) = b(0,1,1)$, whose first two
coordinates force $a = b = 0$): the sum is direct, and $F \oplus
G$ is exactly that plane.
\end{example}

\begin{proposition}\label{prop:b1:vspaces:directsum}
$F + G$ is direct if and only if $F \cap G = \{0\}$.
\end{proposition}

\begin{proof}
If some $w \neq 0$ lies in $F \cap G$: $w = w + 0 = 0 + w$ are two
decompositions of $w$. Conversely, if $u + v = u' + v'$ with $u, u'
\in F$, $v, v' \in G$, then $u - u' = v' - v$ belongs to $F \cap G =
\{0\}$: decompositions are unique.
\end{proof}

\begin{method}[Proving $E = F \oplus G$]
\label{met:b1:vspaces:directsum}
Two things to check, each with its standard opening move.
\begin{enumerate}
  \item \emph{Trivial intersection.} Take $x \in F \cap G$,
        write out both membership conditions, and squeeze $x =
        0$. (Never argue by picture: cf.\ the pitfalls below.)
  \item \emph{Sum is everything.} Take an arbitrary $x \in E$
        and \emph{produce} the decomposition $x = f + g$ ---
        either by guessing $f$ from the target ($f$ must satisfy
        the defining property of $F$, which usually dictates its
        formula) or by solving the linear system expressing $x$
        against generators of $F$ and $G$.
\end{enumerate}
When the decomposition formula is guessed, uniqueness is
automatic from step 1; when only existence is unclear, step 2 is
where the work lives. The two examples below run the method: for
even/odd functions the formula for $f$ is forced by evaluating
the desired identity at $x$ and $-x$; for polynomials vanishing
at a point, by evaluating at $a$.
\end{method}

\begin{example}\label{ex:b1:vspaces:evenodd}
In $\mathcal{F}(\R, \R)$, the even functions $\mathcal{P}$ and the
odd functions $\mathcal{I}$ are supplementary: any $f$ writes
\[
f(x) = \underbrace{\frac{f(x) + f(-x)}{2}}_{\text{even}}
+ \underbrace{\frac{f(x) - f(-x)}{2}}_{\text{odd}},
\]
and a function both even and odd is zero. (Applied to $\exp$, this
is the pair $(\cosh, \sinh)$ of \cref{ch:b1:functions}.)
\end{example}

\begin{example}[{A supplementary pair in $K_n[X]$}]
\label{ex:b1:vspaces:hyperplane}
Fix $a \in K$ and set $F = \{P \in K_n[X] : P(a) = 0\}$, $G =
\operatorname{Vect}(1)$ (the constants). Then $K_n[X] = F \oplus
G$. Indeed $F \cap G$ consists of the constants vanishing at $a$,
i.e.\ $\{0\}$; and every $P$ decomposes as
\[
P = \underbrace{\bigl(P - P(a)\bigr)}_{\in F}
+ \underbrace{P(a)}_{\in G} .
\]
The decomposition is worth memorizing: subtracting the value at a
point is the standard way to project onto ``functions vanishing at
$a$''. Note that $F$ is a large subspace and $G$ a small one; a
supplementary pair need not be balanced in any sense.
\end{example}

\begin{example}[A supplementary subspace is never unique]
\label{ex:b1:vspaces:manysupplements}
In $\R^2$, let $F = \operatorname{Vect}\bigl((1,0)\bigr)$ (the
$x$-axis). Both $G = \operatorname{Vect}\bigl((0,1)\bigr)$ and
$G' = \operatorname{Vect}\bigl((1,1)\bigr)$ are supplementary to
$F$: each meets $F$ only at $0$, and each pair sums to $\R^2$.
The decompositions of a same vector differ:
\[
(2,\ 1.5) = \underbrace{(2, 0)}_{\in F} +
\underbrace{(0, 1.5)}_{\in G}
= \underbrace{(0.5,\ 0)}_{\in F} +
\underbrace{(1.5,\ 1.5)}_{\in G'} .
\]
In fact \emph{every} line other than $F$ itself is a
supplementary of $F$ in $\R^2$: supplementaries are plentiful,
and speaking of ``the'' supplementary is meaningless until an
extra structure (an inner product, \cref{ch:b1:euclid}) singles
one out.
\end{example}

\begin{omfigure}
\begin{tikzpicture}[scale=1.5]
  \draw[->, thin] (-0.4,0) -- (2.9,0) node[below] {$x$};
  \draw[->, thin] (0,-0.4) -- (0,1.9) node[left] {$y$};
  \draw[omDef, very thick] (-0.3,0) -- (2.7,0)
    node[below, pos=0.95, font=\small] {$F$};
  \draw[omThm, very thick] (0,-0.3) -- (0,1.7)
    node[left, pos=0.95, font=\small] {$G$};
  \draw[omProp, very thick] (-0.25,-0.25) -- (1.75,1.75)
    node[right, pos=0.97, font=\small] {$G'$};
  \fill (2,1.5) circle (0.045) node[above right, font=\small]
    {$(2, 1.5)$};
  \draw[omThm, dashed] (2,1.5) -- (2,0);
  \draw[omDef, dashed] (2,1.5) -- (0,1.5);
  \draw[omProp, dashed] (2,1.5) -- (0.5,0);
  \fill (2,0) circle (0.035);
  \fill (0.5,0) circle (0.035);
\end{tikzpicture}

{\small Two decompositions of the same point of $\R^2$ along $F$
(the $x$-axis): with supplementary $G$ (vertical drop) and with
supplementary $G'$ (oblique drop). The $F$-components differ:
a projection depends on the direction of descent.}
\end{omfigure}

\section{Free families, generating families, bases}

\begin{definition}\label{def:b1:vspaces:free}
A family $(x_1, \dots, x_p)$ of vectors of $E$ is:
\begin{itemize}
  \item \emph{generating} (of $E$) when
        $\operatorname{Vect}(x_1,\dots,x_p) = E$;
  \item \emph{free}\index{free family} (its vectors \emph{linearly
        independent}) when
        \[
        \lambda_1 x_1 + \dots + \lambda_p x_p = 0
        \implies \lambda_1 = \dots = \lambda_p = 0 ;
        \]
        otherwise \emph{linked};
  \item a \emph{basis}\index{basis} when it is free and generating.
\end{itemize}
\end{definition}

\begin{proposition}[Coordinates]\label{prop:b1:vspaces:coordinates}
$(e_1, \dots, e_n)$ is a basis of $E$ if and only if every $x \in E$
is \emph{uniquely} a combination $x = \lambda_1 e_1 + \dots +
\lambda_n e_n$; the scalars $\lambda_i$ are the
\emph{coordinates}\index{coordinates} of $x$ in the basis.
\end{proposition}

\begin{proof}
Generating $=$ existence of the decomposition. Uniqueness $=$
freeness: two decompositions of the same $x$ differ by a combination
equal to $0$; freeness forces all its coefficients --- the
differences of coordinates --- to vanish. Conversely, a nontrivial
null combination gives the two decompositions $0 = \sum \lambda_i
e_i = \sum 0\,e_i$.
\end{proof}

\begin{example}\label{ex:b1:vspaces:bases}
The \emph{canonical basis} of $K^n$: $e_i = (0, \dots, 1, \dots, 0)$
($1$ in slot $i$). The monomials $(1, X, X^2, \dots, X^n)$: a basis
of $K_n[X]$ (freeness: a null combination is the zero polynomial, so
all coefficients vanish, \cref{def:b1:poly:def}). In
$\C$ over $\R$: the basis $(1, \iu)$.
\end{example}

\begin{remark}[Coordinates are a team effort]
\label{rem:b1:vspaces:teameffort}
The first coordinate of $x$ in a basis $(e_1, \dots, e_n)$
depends on \emph{all} the basis vectors, not just $e_1$. In
$\R^2$: the vector $(3, 1)$ has first coordinate $3$ in the
canonical basis, but first coordinate $2$ in the basis
$\bigl((1,0), (1,1)\bigr)$ --- solve $(3,1) = a(1,0) + b(1,1)$:
$b = 1$, $a = 2$. Changing one basis vector reshuffles
\emph{every} coordinate; \cref{ch:b1:matrices} will package this
reshuffling into the change-of-basis matrix.
\end{remark}

\begin{example}[Testing a candidate basis, start to finish]
\label{ex:b1:vspaces:testbasis}
Is $\mathcal{F} = (1 + X,\ 1 + X^2,\ X + X^2)$ a basis of
$\R_2[X]$? Write $u_1, u_2, u_3$ for the three polynomials.
\emph{Freeness}: a null combination $a\,u_1 + b\,u_2 + c\,u_3 =
0$ gives, coefficient by coefficient,
\[
a + b = 0, \qquad a + c = 0, \qquad b + c = 0 ;
\]
subtracting the first two, $b = c$, then the third gives $2b =
0$: $a = b = c = 0$, free. \emph{Generating}: instead of solving
three systems, notice the symmetric combination
\[
u_1 + u_2 - u_3 = (1 + X) + (1 + X^2) - (X + X^2) = 2 ,
\]
so $1 = \frac12(u_1 + u_2 - u_3)$; then
\[
X = u_1 - 1 = \tfrac12\bigl(u_1 - u_2 + u_3\bigr),
\qquad
X^2 = u_2 - 1 = \tfrac12\bigl(-u_1 + u_2 + u_3\bigr).
\]
The monomials lie in the span, so everything does: $\mathcal{F}$
is a basis. As a bonus, assembling the three displays gives the
coordinates of any $P = \alpha + \beta X + \gamma X^2$:
\[
P = \frac{\alpha + \beta - \gamma}{2}\,u_1
+ \frac{\alpha - \beta + \gamma}{2}\,u_2
+ \frac{-\alpha + \beta + \gamma}{2}\,u_3 .
\]
(Sanity check with $P = X$: coordinates $\bigl(\frac12,
-\frac12, \frac12\bigr)$, as found above.) Two lessons: symmetry
in the family usually hides a shortcut combination; and once
dimension is available (\cref{ch:b1:findim}), the whole
generating half of this work will come free of charge --- three
free vectors of a $3$-dimensional space always form a basis.
\end{example}

\begin{proposition}[Useful freeness criteria]\label{prop:b1:vspaces:freecriteria}
\begin{enumerate}
  \item A family of \emph{nonzero polynomials of pairwise distinct
        degrees} is free.
  \item Adding a vector to a free family keeps it free if and only
        if the vector is outside the span of the family.
  \item Any subfamily of a free family is free; any family
        containing a generating family is generating.
\end{enumerate}
\end{proposition}

\begin{proof}
(1) In a null combination, look at the highest degree present: its
coefficient must vanish (nothing cancels that degree), and cascade
down.

(2) If $x \in \operatorname{Vect}(x_1, \dots, x_p)$, the relation $x
- \sum\lambda_i x_i = 0$ is nontrivial. Conversely, a nontrivial
null combination of $(x_1, \dots, x_p, x)$ must involve $x$ with a
nonzero coefficient (else it contradicts freeness of the small
family), and solving for $x$ puts it in the span.

(3) Subfamily: a null combination of the subfamily is one of the
whole family with the missing coefficients set to $0$; freeness
of the big family kills all of them. Superfamily: every vector
of $E$ is already a combination of the generating part; give the
extra vectors the coefficient $0$.
\end{proof}

\begin{example}[The staircase principle]
\label{ex:b1:vspaces:staircase}
Let $P_0, P_1, \dots, P_n \in K_n[X]$ with $\deg P_k = k$ for each
$k$ (a ``staircase'' of degrees). Then $(P_0, \dots, P_n)$ is a
basis of $K_n[X]$. Freeness is
\cref{prop:b1:vspaces:freecriteria}~(1). For the generating
property, argue by finite descent on the degree: let $Q \in
K_n[X]$, $Q \neq 0$, of degree $d$, with leading coefficient $a$,
and let $b \neq 0$ be the leading coefficient of $P_d$. Then $Q -
\frac ab P_d$ has degree $< d$ (the top terms cancel); replacing
$Q$ by this difference and iterating, after at most $n + 1$ steps
one reaches the zero polynomial, and unwinding the subtractions
expresses $Q$ as a combination of the $P_k$. Two staircases
already met: the shifted powers $\bigl((X-a)^k\bigr)_{0 \leq k
\leq n}$ (\cref{exo:b1:vspaces:4}), and the Newton products
$\bigl((X - x_0)(X - x_1)\cdots(X - x_{k-1})\bigr)_{0 \leq k \leq
n}$, put to work in the weekend problem.
\end{example}

\begin{example}[Freeness in function spaces]\label{ex:b1:vspaces:functionfree}
In $\mathcal{F}(\R,\R)$, the family $(\eu^{a_1 x}, \dots, \eu^{a_p
x})$ with $a_1 < \dots < a_p$ is free: divide a null combination by
$\eu^{a_p x}$ and let $x \to +\infty$; the last coefficient dies,
and one cascades down (\cref{exo:b1:vspaces:8} details this and
variants). Freeness of functions is proved by \emph{evaluating}:
at well-chosen points, at infinity, or after differentiating.
\end{example}

\begin{example}[A hidden relation shrinks a span]
\label{ex:b1:vspaces:hiddenrelation}
In $\mathcal{F}(\R, \R)$, what is
$\operatorname{Vect}(1,\ \cos^2,\ \sin^2)$? The identity
$\cos^2 + \sin^2 = 1$ is a nontrivial null combination
\[
1\cdot\mathbf{1} + (-1)\cos^2 + (-1)\sin^2 = 0 :
\]
the family is linked, and the span is already generated by
$(1, \cos^2)$ alone ($\sin^2 = 1 - \cos^2$). That smaller family
is free: $a + b\cos^2 x = 0$ for all $x$ gives, at $x = 0$ and
$x = \frac\pi2$: $a + b = 0$ and $a = 0$. So the span is a
\emph{plane} inside the function space --- and it also contains
$\cos 2x = 2\cos^2 x - 1$: linear-looking families of
trigonometric functions routinely collapse under identities,
which is why freeness must be \emph{proved}, never assumed from
the list's length.
\end{example}

\begin{remark}[Common pitfalls]
\label{rem:b1:vspaces:pitfalls}
Four classical traps.
\emph{Pairwise is not enough}: in $\R^2$, the vectors $(1,0)$,
$(0,1)$, $(1,1)$ are pairwise non-proportional, yet linked ---
freeness is a property of the \emph{whole} family, tested by one
global combination, never two by two.
\emph{The zero vector poisons everything}: any family containing
$0$ is linked ($1\cdot 0 = 0$ is a nontrivial relation), however
innocent the other vectors.
\emph{Union is not sum}: $F \cup G$ is almost never a subspace
(\cref{def:b1:vspaces:subspace}); the smallest subspace
containing both is $F + G$, usually much bigger than the union
--- in $\R^2$, two distinct lines have union a cross, sum the
whole plane.
\emph{Direct requires trivial intersection, not disjointness}:
two subspaces are never disjoint (both contain $0$); the correct
condition is $F \cap G = \{0\}$, and it must be \emph{proved},
not read off a drawing --- cf.\
\cref{ex:b1:vspaces:manysupplements}, where many different $G$
work.
\emph{Freeness depends on the scalars}: the pair $(1, \iu)$ is
free in $\C$ seen as an $\R$-vector space, but linked in $\C$
seen as a $\C$-vector space ($\iu\cdot 1 + (-1)\cdot\iu = 0$).
Always know which field is acting before declaring a family free
--- the weekend problem of \cref{ch:b1:findim} turns exactly
this sensitivity into irrationality proofs.
\end{remark}

\begin{remark}[Where this language goes]
\label{rem:b1:vspaces:whereused}
Everything after this chapter speaks the language set up here.
\cref{ch:b1:findim} counts basis vectors and turns ``free'' and
``generating'' into inequalities on a single integer, the
dimension. \cref{ch:b1:linmaps} studies the maps compatible with
the two operations; direct sums become projectors there.
\cref{ch:b1:matrices} encodes vectors by their coordinates in a
basis --- \cref{prop:b1:vspaces:coordinates} is the licence for
that encoding --- and \cref{ch:b1:euclid} adds lengths and angles
on top of the linear structure. In the Year~2 volume the same
axioms, verbatim, run over arbitrary fields and in infinite
dimension; nothing in this chapter used finiteness anywhere.
\end{remark}

\begin{remark}[Three threads to follow through Book 3]
\label{rem:b1:vspaces:threads}
Watch three specific ideas of this chapter grow.
\emph{The staircase principle}
(\cref{ex:b1:vspaces:staircase}) reappears as the Newton basis
in this chapter's weekend problem, as the binomial basis
$(B_k)$ there, and as the polynomial alternant trick in the
weekend problem of \cref{ch:b1:det}: one lemma, three
determinant-free dividends.
\emph{Evaluation as a freeness test}
(\cref{ex:b1:vspaces:functionfree}) becomes the interpolation
isomorphism of \cref{ch:b1:linmaps}, then the Vandermonde
criterion of \cref{ch:b1:det}, then the Gram test of
\cref{ch:b1:euclid}: the same reflex, sharpened three times.
\emph{Direct sums} (\cref{def:b1:vspaces:sum}) become
projectors in \cref{ch:b1:linmaps}, orthogonal splittings $E =
F \oplus F^\perp$ in \cref{ch:b1:euclid}, and the
explained-plus-residual decomposition of least squares in the
weekend problem of \cref{ch:b1:multivar}. Very little of this
book is not, at bottom, one of these three ideas wearing new
clothes.
\end{remark}

\section{Exercises}

\begin{exercise}[$\star$]\label{exo:b1:vspaces:1}
Which of the following are subspaces?
\begin{enumerate}
  \item $\{(x, y, z) \in \R^3 : x + 2y - z = 0\}$;
  \item $\{(x, y, z) \in \R^3 : x + 2y - z = 1\}$;
  \item $\{(x, y) \in \R^2 : xy \geq 0\}$;
  \item $\{P \in \R[X] : P(1) = 0\}$;
  \item $\{f \in \mathcal{F}(\R,\R) : f \text{ bounded}\}$.
\end{enumerate}
\end{exercise}

\begin{exercise}[$\star$]\label{exo:b1:vspaces:2}
In $\R^3$, is $(1, 2, 1)$ in
$\operatorname{Vect}\bigl((1,0,1),\, (1,1,0)\bigr)$? And $(2, 1,
1)$? Describe $\operatorname{Vect}\bigl((1,0,1),(1,1,0)\bigr)$ by an
equation.
\end{exercise}

\begin{exercise}[$\star$]\label{exo:b1:vspaces:3}
Decide freeness in $\R^3$: $\;\bigl((1,1,0), (1,0,1),
(0,1,1)\bigr)$; $\;\bigl((1,2,3), (2,4,6)\bigr)$;
$\;\bigl((1,0,0), (1,1,0), (1,1,1), (0,1,1)\bigr)$.
\end{exercise}

\begin{exercise}[$\star$]\label{exo:b1:vspaces:4}
Prove that $(1, X - 1, (X-1)^2, (X-1)^3)$ is a basis of $\R_3[X]$,
and give the coordinates of $X^3$ in it. \emph{(Taylor at $1$!)}
\end{exercise}

\begin{exercise}[$\star\star$]\label{exo:b1:vspaces:5}
In $\R^4$, let $F = \{(x,y,z,t) : x = y = z\}$ and $G = \{(x,y,z,t)
: x = t = 0\}$. Prove that $F \oplus G = \R^4$, and decompose
$(1,2,3,4)$ accordingly.
\end{exercise}

\begin{exercise}[$\star\star$]\label{exo:b1:vspaces:6}
In the space of sequences, let $F$ be the set of convergent
sequences and $G = \operatorname{Vect}(u)$ where $u_n = (-1)^n$.
Prove that $F \cap G = \{0\}$. Is $F + G$ the whole space of
sequences?
\end{exercise}

\begin{exercise}[$\star\star$]\label{exo:b1:vspaces:7}
Let $F, G, H$ be subspaces of $E$. Prove that
\[
F \cap (G + (F \cap H)) = (F \cap G) + (F \cap H),
\]
and show by an example in $\R^2$ that the unrestricted
distributivity $F \cap (G + H) = (F\cap G) + (F \cap H)$ fails.
\end{exercise}

\begin{exercise}[$\star\star\star$]\label{exo:b1:vspaces:8}
Prove that the following families of $\mathcal{F}(\R, \R)$ are free:
\begin{enumerate}
  \item $(\eu^{a_1 x}, \dots, \eu^{a_p x})$ for $a_1 < \dots < a_p$;
  \item $(\cos x, \sin x, \cos 2x, \sin 2x)$;
  \item $(x \mapsto \abs{x - a_1}, \dots, x \mapsto \abs{x - a_p})$
        for distinct $a_i$ \emph{(differentiability fails at exactly
        one point per function)}.
\end{enumerate}
\end{exercise}

\begin{exercise}[$\star\star\star$]\label{exo:b1:vspaces:9}
Let $E$ be a $K$-vector space and $F, G, H$ subspaces with $F + G =
F + H$, $F \cap G = F \cap H$ and $G \subseteq H$. Prove $G = H$.
Give a counterexample without the hypothesis $G \subseteq H$.
\end{exercise}

\begin{exercise}[$\star\star$]\label{exo:b1:vspaces:10}
In $\R[X]$, let $\mathcal P$ be the set of \emph{even} polynomials
($P(-X) = P(X)$) and $\mathcal I$ the set of \emph{odd} ones
($P(-X) = -P(X)$). Prove that $\R[X] = \mathcal P \oplus \mathcal
I$, and show that $\mathcal P =
\operatorname{Vect}(1, X^2, X^4, \dots)$, i.e.\ that the even
polynomials are exactly the polynomials in $X^2$.
\end{exercise}

\begin{exercise}[$\star\star$]\label{exo:b1:vspaces:11}
Let $(x_1, x_2, x_3)$ be a free family of a real vector space $E$.
Prove that $(x_1 + x_2,\ x_2 + x_3,\ x_3 + x_1)$ is free. Is the
analogous family of four vectors $(x_1 + x_2,\ x_2 + x_3,\ x_3 +
x_4,\ x_4 + x_1)$ free when $(x_1, x_2, x_3, x_4)$ is?
\end{exercise}

\begin{exercise}[$\star\star\star$]\label{exo:b1:vspaces:12}
Let $E$ be a vector space over $\R$ (or $\C$) and $F_1, \dots,
F_k$ \emph{proper} subspaces of $E$ (each $F_i \neq E$).
\begin{enumerate}
  \item Treat the case $k = 2$ directly: if $F_1 \not\subseteq
        F_2$ and $F_2 \not\subseteq F_1$, pick $x \in F_1
        \setminus F_2$ and $y \in F_2 \setminus F_1$ and locate $x
        + y$.
  \item Prove in general that $E \neq F_1 \cup \dots \cup F_k$: a
        vector space over an infinite field is never a finite
        union of proper subspaces. \emph{(Take $k$ minimal, pick
        $x \in F_1$ outside the other $F_i$, pick $y \notin F_1$,
        and follow the line $t \mapsto y + tx$.)}
\end{enumerate}
\end{exercise}

\section{Problem: interpolation, three bases for one space}

\begin{problem}\label{pb:b1:vspaces:1}
Fix $n + 1$ \emph{distinct} points $x_0, x_1, \dots, x_n$ of $\R$.
This problem revisits Lagrange interpolation
(\cref{thm:b1:poly:lagrange}) with the eyes of this chapter: the
space $\R_n[X]$ carries three natural bases --- Lagrange's,
Newton's, and (for equally spaced points) the binomial basis ---
and each basis makes one question easy. The road ends at a genuine
arithmetic theorem: P\'olya's characterization of the polynomials
mapping $\Z$ into $\Z$.\index{Lagrange interpolation}
\index{divided differences}\index{integer-valued polynomial}

\medskip
\noindent\textbf{Part I --- The Lagrange basis.} For $0 \leq i
\leq n$ set
\[
L_i \;=\; \prod_{j \neq i} \frac{X - x_j}{x_i - x_j} \;\in\;
\R_n[X].
\]
\begin{enumerate}
  \item Check that $\deg L_i = n$ and that $L_i(x_j) = 1$ if $j =
        i$, and $0$ if $j \neq i$.
  \item Prove that the family $(L_0, \dots, L_n)$ is free.
  \item Prove that for every $P \in \R_n[X]$,
        \[
        P \;=\; \sum_{i=0}^{n} P(x_i)\, L_i ,
        \]
        and conclude that $(L_0, \dots, L_n)$ is a basis of
        $\R_n[X]$. \emph{(Consider the difference of the two sides
        and count its roots, \cref{cor:b1:poly:nroots}.)}
  \item Deduce the interpolation theorem: for any values $y_0,
        \dots, y_n \in \R$ there is a \emph{unique} $P \in
        \R_n[X]$ with $P(x_i) = y_i$ for all $i$. In the Lagrange
        basis, what are the coordinates of a polynomial $P$?
  \item Prove the identities
        \[
        \sum_{i=0}^{n} L_i = 1
        \qquad\text{and, for } 0 \leq k \leq n,\qquad
        \sum_{i=0}^{n} x_i^{k}\, L_i = X^{k} .
        \]
\end{enumerate}

\medskip
\noindent\textbf{Part II --- The Newton basis and divided
differences.} Set $N_0 = 1$ and $N_k = (X - x_0)(X - x_1) \cdots
(X - x_{k-1})$ for $1 \leq k \leq n$. For a function $f$ defined
at the nodes, define the \emph{divided differences} by
$f[x_i] = f(x_i)$ and
\[
f[x_i, \dots, x_{i+k}] \;=\;
\frac{f[x_{i+1}, \dots, x_{i+k}] - f[x_i, \dots, x_{i+k-1}]}
{x_{i+k} - x_i} .
\]
\begin{enumerate}[resume]
  \item Prove that $(N_0, N_1, \dots, N_n)$ is a basis of
        $\R_n[X]$.
  \item Compute $f[x_0, x_1]$ and $f[x_0, x_1, x_2]$ in terms of
        the values of $f$, then compute all divided differences of
        $f(x) = x^2$ at three arbitrary nodes.
  \item (Aitken's lemma) Let $R$ interpolate $f$ at $x_0, \dots,
        x_{n-1}$ and $Q$ interpolate $f$ at $x_1, \dots, x_n$,
        both of degree $\leq n - 1$. Prove that
        \[
        S \;=\; \frac{(X - x_0)\,Q - (X - x_n)\,R}{x_n - x_0}
        \]
        interpolates $f$ at $x_0, x_1, \dots, x_n$.
  \item Deduce, by induction on the number of nodes, that the
        coefficient of $X^{k}$ in the interpolant of $f$ at $x_0,
        \dots, x_k$ is exactly $f[x_0, \dots, x_k]$.
  \item Prove \emph{Newton's interpolation formula}: the
        interpolant of $f$ at $x_0, \dots, x_n$ is
        \[
        P \;=\; \sum_{k=0}^{n} f[x_0, \dots, x_k]\, N_k ,
        \]
        and derive the closed formula
        \[
        f[x_0, \dots, x_k] \;=\; \sum_{i=0}^{k}
        \frac{f(x_i)}{\prod_{j \neq i,\, j \leq k} (x_i - x_j)} ,
        \]
        which shows that $f[x_0, \dots, x_k]$ does not depend on
        the ordering of the nodes.
\end{enumerate}

\medskip
\noindent\textbf{Part III --- Equally spaced nodes: the difference
operator.} From now on the nodes are $0, 1, 2, \dots$ and, for a
polynomial $P$, we set
\[
\Delta P(X) = P(X + 1) - P(X),
\qquad
B_k = \frac{X(X-1)\cdots(X-k+1)}{k!} \quad (B_0 = 1).
\]
\begin{enumerate}[resume]
  \item Show that if $\deg P = m \geq 1$ with leading coefficient
        $a$, then $\deg \Delta P = m - 1$ with leading coefficient
        $m\,a$, and that $\Delta$ kills constants.
  \item Show that $(B_0, B_1, \dots, B_n)$ is a basis of
        $\R_n[X]$ and that $\Delta B_k = B_{k-1}$ for $k \geq 1$.
  \item (Newton's forward-difference formula) Prove that every $P
        \in \R_n[X]$ satisfies
        \[
        P \;=\; \sum_{k=0}^{n} \bigl(\Delta^{k} P\bigr)(0)\, B_k .
        \]
  \item Prove that for every $k \geq 0$,
        \[
        \bigl(\Delta^{k} P\bigr)(0) \;=\;
        \sum_{j=0}^{k} (-1)^{k-j} \binom{k}{j} P(j) .
        \]
  \item Show that if $\deg P = n$ with leading coefficient $a_n$,
        then $\Delta^{n} P$ is the constant $n!\,a_n$ and
        $\Delta^{n+1} P = 0$.
\end{enumerate}

\medskip
\noindent\textbf{Part IV --- Integer-valued polynomials.} A
polynomial $P \in \R[X]$ is \emph{integer-valued} when $P(m) \in
\Z$ for every $m \in \Z$.
\begin{enumerate}[resume]
  \item Prove that each $B_k$ is integer-valued. \emph{(Treat $m
        \geq k$, $0 \leq m < k$ and $m < 0$ separately; for $m =
        -q < 0$, show $B_k(-q) = (-1)^k \binom{q + k - 1}{k}$.)}
  \item Prove \emph{P\'olya's characterization}: $P \in \R_n[X]$
        is integer-valued if and only if its coordinates in the
        basis $(B_0, \dots, B_n)$ are integers.
  \item Deduce: if $P \in \R_n[X]$ takes integer values at $n + 1$
        \emph{consecutive} integers $a, a+1, \dots, a+n$, then $P$
        is integer-valued. \emph{(Shift: apply the study to $Q(X)
        = P(X + a)$.)}
  \item Deduce from question 16 that a product of $k$ consecutive
        integers is always divisible by $k!$.
  \item Let $P = \dfrac{X(X+1)(2X+1)}{6}$. Compute its Newton
        table at $0, 1, 2, 3$, write $P$ in the basis $(B_k)$, and
        conclude that $P$ is integer-valued although none of its
        monomial coefficients is an integer. Verify $\Delta P =
        (X+1)^2$ and deduce $P(m) = 1^2 + 2^2 + \dots + m^2$ for
        $m \in \N$.
\end{enumerate}

\medskip
\noindent\textbf{Part V --- Dividends.}
\begin{enumerate}[resume]
  \item Let $P \in \R_n[X]$ interpolate the values $2^i$ at $i =
        0, 1, \dots, n$. Show that $P = B_0 + B_1 + \dots + B_n$
        and that $P(n + 1) = 2^{n+1} - 1$: the ``doubling
        pattern'' always breaks at the very next point.
  \item (Discrete antiderivative) Prove that for all integers $m
        \geq 1$ and $k \geq 0$,
        \[
        \sum_{j=0}^{m-1} B_k(j) \;=\; B_{k+1}(m),
        \]
        i.e.\ the hockey-stick identity $\sum_{j=k}^{m-1}
        \binom{j}{k} = \binom{m}{k+1}$.
  \item Expand $X^2$ and $X^3$ in the basis $(B_k)$ and deduce
        closed formulas for $\sum_{j=0}^{m-1} j^2$ and
        $\sum_{j=0}^{m-1} j^3$; recover Nicomachus' identity
        $1^3 + \dots + m^3 = (1 + \dots + m)^2$.
  \item Take $n = 2$ and nodes $0, 1, 2$. Write the coordinates of
        $X^2$ in the three bases of this problem: the monomial
        basis, the Lagrange basis, the Newton basis. Check the
        three answers against questions 4 and 9.
  \item Synthesis. In four sentences: which vector-space concept
        makes question 4 automatic; why the Newton basis computes
        coordinates \emph{recursively} while the Lagrange basis
        reads them off \emph{instantly}; which freeness criterion
        both bases share; and in what precise sense P\'olya's
        theorem says that integrality of a polynomial is a
        property of its coordinates \emph{in the right basis}.
\end{enumerate}
\end{problem}
